\documentclass[10pt,a4paper]{amsart}
\usepackage[margin=19mm]{geometry}
\usepackage{amsmath,amssymb,mathtools}
\usepackage[scr=rsfs,cal=euler]{mathalfa}
\usepackage{xfrac}
\usepackage[unicode,hidelinks,bookmarks=false]{hyperref}
\newcommand{\ie}{i.e.\ }
\newcommand{\cf}{cf.\ }
\newcommand{\resp}{respectively\ }
\newcommand{\Subsection}{\subsection}
\providecommand{\nobreakdash}{\nobreak\hbox{-}}
\setlength{\emergencystretch}{2em}
\multlinegap0pt
\usepackage{xcolor}
\usepackage[all]{xy}
\usepackage{amssymb}
\newtheorem{lemma}{Lemma}
\renewcommand{\H}{\mathbb{H}}
\renewcommand{\P}{\mathbb{P}}
\newcommand{\R}{\mathbb{R}}
\newcommand{\Z}{\mathbb{Z}}
\title{Valuations on the character variety: \\ Newton polytopes and Residual Poisson Bracket}
\author{Julien March\'e, Christopher-Lloyd Simon}
\date{Source: arXiv:2104.04340v1 (CC BY 4.0)}
\begin{document}
\maketitle

\noindent\emph{Source and licence.} Valuations on the character variety: \\ Newton polytopes and Residual Poisson Bracket. Version arXiv:2104.04340v1 (CC BY 4.0), distributed by arXiv under the Creative Commons Attribution 4.0 International licence, http://creativecommons.org/licenses/by/4.0/, as recorded in the arXiv licence metadata for this exact version and shown on the abstract page https://arxiv.org/abs/2104.04340v1. Full licence notice: \texttt{licenses/arxiv-2104.04340-CC-BY-4.0.txt}.\par\smallskip

\noindent\textit{Editorial scope.} Counted unit: the article's lemma equiv-homotopy (source0.tex lines 1192--1194). The article's complete original proof of it is delivered with it (source0.tex lines 1195--1241, one \texttt{proof} environment of 47 lines). No mathematics is written, repaired or re-derived: the statement and the proof are the article's own lines, in the article's own notation, and no retained line is substituted or re-flowed.  Retained uncounted as the source-local context the proof needs: lines 1249--1251.  One declared adaptation: exactly 1 ancillary strings inside the retained spans are omitted (image-inclusion commands and, where the source repeats a label, the repeated label definitions) and no other line differs from the source; the figure environments, with their placement, captions and labels, are kept, and the package records the omitted strings in fidelity receipts bound to the affected bodies.  No retained line contains a citation, so the wrapper carries no bibliography and no bibliography entry.  The retained lines contain the article's own diagram code, which the wrapper typesets with the standard tikz packages; no image file is delivered and the package declares no asset dependency.  Editorial formatting only, in addition: an AMS \texttt{amsart} preamble that restates the article's own declarations and macros used by the retained lines, namely the environments \texttt{lemma} on separate counters, together with the standard packages those lines need.  The retained environments print in this document as Lemma 1 in order of appearance, whereas the article prints the same items with the numbering of its own full document.  The complete native source of the article as distributed in the arXiv e-print for this exact version is bundled unchanged as \texttt{sources/110108/source0.tex}.
\begin{lemma}\label{retraction}
For all finite $W\subset V(T)$, the cell-complex $T^o_{(W)}$ retracts by deformation on $T^o_W$ .
\end{lemma}

\begin{lemma}\label{equiv-homotopy}
Let $\tilde{S}^o$ be the covering of the orbifold $S^o$ corresponding to the kernel of the natural map $\pi_1(S^o)\to \pi_1(S)$. There exists a $\pi_1(S)$-equivariant map $F:T^o\to \tilde{S}^o$ which induces a homotopy equivalence between $\Sigma^o$ and $S^o$.
\end{lemma}

\begin{proof}


To define $F$, represent $T$ as the dual tree to a measured geodesic lamination $\lambda$ and consider the collection of circles inscribed in each triangle of the complement $S\setminus\lambda$: they lift to a collection of circles $C_x$ in $\tilde{S}\simeq \H^2$ indexed by $x\in V(T)$.
%
Moreover, the half-edges incident to $x$ correspond bijectively to the three intersection points of $C_x$ with the leaves of the lamination, see Figure \ref{fig:equivmap}.

\begin{figure}[htbp]
    \centering
    
    \caption{Lifting a geodesic to $\H^2$ (done with Geogebra)}
    \label{fig:equivmap}
\end{figure}
%
%(Notice that since $\lambda$ is filling, the discs interior to the $C_x$ have disjoint neighborhoods, and the minimal hyperbolic distance between two of them can be read in $S$.)
%
The covering $\tilde{S}^o$ is obtained from $\H^2$ by drilling out the interior of $C_x$ and gluing back a copy of $\R\P^\infty$ along $\R\P^1\simeq C_x$ for every $x\in V(T)$. By construction, the orbifold $S^o$ is homotopic to the quotient $\tilde{S}^o/\pi_1(S)$. 

We now proceed to the construction of an equivariant map $F \colon T^o\to \tilde{S}^o$. There is already an identification between the orbifold parts of both spaces, so that we are left to define the map $F$ on the connecting part. 

For every pair $(x,y)\in V(T)^2$, we must define a path $F(x,y)$ in $\tilde{S}^o$ connecting the points of $C_x$ and $C_y$ identified to the end points $h_x,h_y$ of $(x,y)$ in $T^o$. A first guess would be to consider the geodesic path $\gamma$ between the points $h_x$ and $h_y$. This path actually projects to the geodesic joining $x$ to $y$ in $T$. However it may intersect a forbidden circle $C_z$, in which case it enters its circumscribed ideal triangle $\Delta_z$ by one side and leaves it by another. Call $p_z$ the ideal vertex at the intersection of these two sides. 
We can homotope $\gamma$ inside $\Delta_z$ to a path avoiding $C_z$ which stays on the side containing $p_z$, see Figure \ref{fig:equivmap}.

%\textcolor{red}{Pourquoi de ce côté et pas de l'autre ? Je ne penses pas qu'on ait le choix sur le côté, et d'ailleurs si on fait certaines alternances gauche droite entre les triangles il faut couper le cercle inscrit...}
%
%This path will indeed belong to $f^{-1}([x,y])$ where $[x,y]$ denotes the geodesic joining $x$ to $y$ and will avoid the inscribed circles included in that region. 
%
%Let $U$ be the set of vertices lying between $x$ and $y$. These vertices belong to two disjoint classes $U=U_L\cup U_R$: a branch point $u$ belongs to $U_L$ (resp. $U_R$) if the branch leaving $u$ turns left when moving from $x$ to $y$. We define $F(x,y)$ to be any transverse path in $f^{-1}([x,y])$ which separate the collection of circles $U$ accordingly to this decomposition. 

Moreover, we can choose those paths in such a way that $F$ is $\pi_1(S)$-equivariant. Let us now consider a triple of points $x,z,y$ lying on a geodesic of $T$ in that order. We have defined $F(x,z)$, $F(z,y)$ and $F(x,y)$: it is not hard to see that the region enclosed by the three arcs and the boundary of $C_z$ does not contain any other circle, hence it can be filled by a triangle: this extends $F$ to the $2$-skeleton of $T^o$. 
%
This procedure can be continued to define an equivariant map $F:T^o\to \tilde{S}^o$, which induces a map $\overline{F}:\Sigma^o\to S^o$.

We would like to show that $\overline{F}$ is a homotopy equivalence. The space $\tilde{S}^o$ is Eilenberg-MacLane and Lemma \ref{retraction} below shows that so is $T^o$, hence it is sufficient to prove that $\overline{F}$ induces an isomorphism between fundamental groups.
%
Behold the following commutative diagram, and observe that the five lemma reduces the statement to showing that $F_*$ is an isomorphism.
$$\xymatrix{
0\ar[r]& \pi_1(T^o)\ar[d]^{F_*}\ar[r] & \pi_1(\Sigma^o)\ar[r]\ar[d]^{\overline{F}_*}& \pi_1(S)\ar[d]^=\ar[r]& 0 \\
0\ar[r]& \pi_1(\tilde{S}^o)\ar[r] & \pi_1(S^o)\ar[r]& \pi_1(S)\ar[r]& 0 
}$$
This last statement is clear from the fact that $\pi_1(T^o)$ and $\pi_1(\tilde{S}^o)$ are both isomorphic to a free product of copies of $\Z/2\Z$ indexed by $V(T)$ (see again Lemma \ref{retraction}). 
%As $T^o$ is an Eilengerg-Maclane space, the same is true for $\Sigma^o$ hence it is sufficient to show that $\pi_1(\Sigma^o)$ and $\pi_1(S^o)$ are isomorphic groups. 
%Fix a hyperbolic structure on $S$ and represent $T$ as the dual tree of a measured lamination $\lambda$ on $S$. Denote by $f$ the quotient map $f:\tilde{S}\simeq \H^2\to T$. 
%Denote by $\tilde{S}^o$ the covering of $S^o$ corresponding to the morphism $\pi_1(S^o)\to \pi_1(S)$. 
%It is sufficient to show that there is an isomorphism $\pi_1(\tilde{S}^o\simeq \pi_1(T^o)$ which is $\pi_1(S)$-equivariant up to conjugation (détailler...). 
%For any finite subset $W\in V(T)$, we can embedd the tree $T'_W$ inside $\H_2$ by sending each cycle to the corresponding inscribed circle. It suffices to lift the edges between closest neighbours to arcs transversal to the lamination...
\end{proof}

\end{document}
